% GENERATED FILE. Edit canonical lessons, then run npm run generate:books.
% canonical-source-hash: fnv1a64:cd3f30bd36f25788
% canonical-lessons: SA-C301-meghah, SA-C301-cakravatah, SA-C301-ghurnate, SA-C301-sphurati, SA-C301-unmulayati, SA-R301-review-294-301

\chapter{The Storm}
\label{ch:sa-the-storm}
\hlchaptermodality{301}

\begin{chapteropening}
\textbf{By the end of this chapter:} \emph{I can describe a storm as it came.}

Complete the last lesson of chapter 301: I can describe a storm as it came.
\end{chapteropening}

\section[\texorpdfstring{\sk{मेघः} (meghaḥ)}{मेघः (meghaḥ)}]{\sk{मेघः} (meghaḥ) — a cloud}

\label{lesson:SA-C301-meghah}



\begin{quote}
Before the new one: say the Sanskrit for cracks open, then the Sanskrit for a pit.
\end{quote}

\subsection*{You'll want to know: \sk{मेघः}}
\textbf{\sk{मेघः}} — \emph{meghaḥ} — \textquotedblleft{}a cloud\textquotedblright{}.

\begin{quote}
\sk{आकाशे} \sk{कृष्णाः} \sk{मेघाः} \sk{सन्ति।}
\end{quote}

\begin{quote}
\emph{ākāśe kṛṣṇāḥ meghāḥ santi.}
\end{quote}

\begin{quote}
There are dark clouds in the sky.
\end{quote}

A masculine noun.

\begin{cousinweb}
\textbf{\sk{मेघः}} (\emph{meghaḥ}) comes from the root √mih (MW writes it √migh here), \textquotedblleft{}to sprinkle\textquotedblright{}: MW explains \emph{megha} as the sprinkler, a cloud (MW).
\end{cousinweb}

\begin{grammarlens}[title={What you've built}]
One more word for telling what happened.
\end{grammarlens}

\subsection*{Your turn}
\begin{itemize}
\raggedright
  \item \emph{Say it:} \emph{meghaḥ}
  \item \emph{Say it:} \emph{meghaḥ}, once more
  \item \emph{Say it:} \emph{ākāśe kṛṣṇāḥ meghāḥ santi.} — There are dark clouds in the sky.
  \item \emph{Recall:} say the Sanskrit for cracks open, then the Sanskrit for a pit, then say \emph{meghaḥ} again
\end{itemize}

\begin{tcolorbox}[breakable,colback=teal!4,colframe=teal!35!black,title={Before you move on}]
What does \sk{मेघः} mean? (\textquotedblleft{}A cloud\textquotedblright{}.) Say it once more.
\end{tcolorbox}
\section[\texorpdfstring{\sk{चक्रवातः} (cakravātaḥ)}{चक्रवातः (cakravātaḥ)}]{\sk{चक्रवातः} (cakravātaḥ) — a whirlwind}

\label{lesson:SA-C301-cakravatah}



\begin{quote}
Before the new one: say the Sanskrit for a pit, then the Sanskrit for a cloud.
\end{quote}

\subsection*{You'll want to know: \sk{चक्रवातः}}
\textbf{\sk{चक्रवातः}} — \emph{cakravātaḥ} — \textquotedblleft{}a whirlwind\textquotedblright{}.

\begin{quote}
\sk{समुद्रात्} \sk{चक्रवातः} \sk{आगच्छति।}
\end{quote}

\begin{quote}
\emph{samudrāt cakravātaḥ āgacchati.}
\end{quote}

\begin{quote}
A whirlwind is coming from the sea.
\end{quote}

A masculine noun.

\begin{cousinweb}
\textbf{\sk{चक्रवातः}} (\emph{cakravātaḥ}) is built from \emph{cakra}, \textquotedblleft{}a wheel\textquotedblright{}, and \emph{vāta}, \textquotedblleft{}wind\textquotedblright{}; MW gives \emph{cakravāta} as a whirlwind, and a cyclone is its modern use (MW).
\end{cousinweb}

\begin{grammarlens}[title={What you've built}]
One more word for telling what happened.
\end{grammarlens}

\subsection*{Your turn}
\begin{itemize}
\raggedright
  \item \emph{Say it:} \emph{cakravātaḥ}
  \item \emph{Say it:} \emph{cakravātaḥ}, once more
  \item \emph{Say it:} \emph{samudrāt cakravātaḥ āgacchati.} — A whirlwind is coming from the sea.
  \item \emph{Recall:} say the Sanskrit for a pit, then the Sanskrit for a cloud, then say \emph{cakravātaḥ} again
\end{itemize}

\begin{tcolorbox}[breakable,colback=teal!4,colframe=teal!35!black,title={Before you move on}]
What does \sk{चक्रवातः} mean? (\textquotedblleft{}A whirlwind\textquotedblright{}.) Say it once more.
\end{tcolorbox}
\section[\texorpdfstring{\sk{घूर्णते} (ghūrṇate)}{घूर्णते (ghūrṇate)}]{\sk{घूर्णते} (ghūrṇate) — whirls; spins round}

\label{lesson:SA-C301-ghurnate}



\begin{quote}
Before the new one: say the Sanskrit for a cloud, then the Sanskrit for a whirlwind.
\end{quote}

\subsection*{You'll want to know: \sk{घूर्णते}}
\textbf{\sk{घूर्णते}} — \emph{ghūrṇate} — \textquotedblleft{}whirls; spins round\textquotedblright{}.

\begin{quote}
\sk{चक्रवातः} \sk{वेगेन} \sk{घूर्णते।}
\end{quote}

\begin{quote}
\emph{cakravātaḥ vegena ghūrṇate.}
\end{quote}

\begin{quote}
The whirlwind spins fast.
\end{quote}

Past (it …): \sk{वायुः} \sk{वेगेन} \sk{घूर्णितवान्}, \emph{vāyuḥ vegena ghūrṇitavān}.

\begin{cousinweb}
\textbf{\sk{घूर्णते}} (\emph{ghūrṇate}) comes from the root √ghūrṇ, \textquotedblleft{}to move to and fro, roll about\textquotedblright{} (MW).
\end{cousinweb}

\begin{grammarlens}[title={What you've built}]
One more word for telling what happened.
\end{grammarlens}

\subsection*{Your turn}
\begin{itemize}
\raggedright
  \item \emph{Say it:} \emph{ghūrṇate}
  \item \emph{Say it:} \emph{ghūrṇate}, once more
  \item \emph{Say it:} \emph{cakravātaḥ vegena ghūrṇate.} — The whirlwind spins fast.
  \item \emph{Recall:} say the Sanskrit for a cloud, then the Sanskrit for a whirlwind, then say \emph{ghūrṇate} again
\end{itemize}

\begin{tcolorbox}[breakable,colback=teal!4,colframe=teal!35!black,title={Before you move on}]
What does \sk{घूर्णते} mean? (\textquotedblleft{}Whirls; spins round\textquotedblright{}.) Say it once more.
\end{tcolorbox}
\section[\texorpdfstring{\sk{स्फुरति} (sphurati)}{स्फुरति (sphurati)}]{\sk{स्फुरति} (sphurati) — flashes; flickers}

\label{lesson:SA-C301-sphurati}



\begin{quote}
Before the new one: say the Sanskrit for a whirlwind, then the Sanskrit for whirls.
\end{quote}

\subsection*{You'll want to know: \sk{स्फुरति}}
\textbf{\sk{स्फुरति}} — \emph{sphurati} — \textquotedblleft{}flashes; flickers\textquotedblright{}.

\begin{quote}
\sk{मेघेषु} \sk{विद्युत्} \sk{स्फुरति।}
\end{quote}

\begin{quote}
\emph{megheṣu vidyut sphurati.}
\end{quote}

\begin{quote}
Lightning flashes in the clouds.
\end{quote}

Past (it …): \sk{विद्युत्} \sk{स्फुरितवती}, \emph{vidyut sphuritavatī}.

\begin{cousinweb}
\textbf{\sk{स्फुरति}} (\emph{sphurati}) comes from the root √sphur, \textquotedblleft{}to dart, quiver, flash\textquotedblright{}; MW compares English \emph{spur} and \emph{spurn} (MW).
\end{cousinweb}

\begin{grammarlens}[title={What you've built}]
One more word for telling what happened.
\end{grammarlens}

\subsection*{Your turn}
\begin{itemize}
\raggedright
  \item \emph{Say it:} \emph{sphurati}
  \item \emph{Say it:} \emph{sphurati}, once more
  \item \emph{Say it:} \emph{megheṣu vidyut sphurati.} — Lightning flashes in the clouds.
  \item \emph{Recall:} say the Sanskrit for a whirlwind, then the Sanskrit for whirls, then say \emph{sphurati} again
\end{itemize}

\begin{tcolorbox}[breakable,colback=teal!4,colframe=teal!35!black,title={Before you move on}]
What does \sk{स्फुरति} mean? (\textquotedblleft{}Flashes; flickers\textquotedblright{}.) Say it once more.
\end{tcolorbox}
\section[\texorpdfstring{\sk{उन्मूलयति} (unmūlayati)}{उन्मूलयति (unmūlayati)}]{\sk{उन्मूलयति} (unmūlayati) — uproots}

\label{lesson:SA-C301-unmulayati}



\begin{quote}
Before the new one: say the Sanskrit for whirls, then the Sanskrit for flashes.
\end{quote}

\subsection*{You'll want to know: \sk{उन्मूलयति}}
\textbf{\sk{उन्मूलयति}} — \emph{unmūlayati} — \textquotedblleft{}uproots\textquotedblright{}.

\begin{quote}
\sk{चक्रवातः} \sk{वृक्षान्} \sk{उन्मूलयति।}
\end{quote}

\begin{quote}
\emph{cakravātaḥ vṛkṣān unmūlayati.}
\end{quote}

\begin{quote}
The whirlwind uproots trees.
\end{quote}

Past (it …): \sk{वायुः} \sk{वृक्षान्} \sk{उन्मूलितवान्}, \emph{vāyuḥ vṛkṣān unmūlitavān}.

\begin{cousinweb}
\textbf{\sk{उन्मूलयति}} (\emph{unmūlayati}) is built from \emph{ud}, \textquotedblleft{}out\textquotedblright{}, and \emph{mūla}, \textquotedblleft{}a root\textquotedblright{}; MW gives \emph{unmūlayati} as to pull up by the roots (MW).
\end{cousinweb}

\begin{grammarlens}[title={What you've built}]
One more word for telling what happened.
\end{grammarlens}

\subsection*{Your turn}
\begin{itemize}
\raggedright
  \item \emph{Say it:} \emph{unmūlayati}
  \item \emph{Say it:} \emph{unmūlayati}, once more
  \item \emph{Say it:} \emph{cakravātaḥ vṛkṣān unmūlayati.} — The whirlwind uproots trees.
  \item \emph{Recall:} say the Sanskrit for whirls, then the Sanskrit for flashes, then say \emph{unmūlayati} again
\end{itemize}

\begin{tcolorbox}[breakable,colback=teal!4,colframe=teal!35!black,title={Before you move on}]
What does \sk{उन्मूलयति} mean? (\textquotedblleft{}Uproots\textquotedblright{}.) Say it once more.
\end{tcolorbox}
\section[(dialogue)]{Looking back: from A Crash on the Road to The Storm}

\label{lesson:SA-R301-review-294-301}



\begin{quote}
Say the Sanskrit for a collision and a blow.
\end{quote}

\subsection*{Your turn}
Say these aloud:

\begin{itemize}
\raggedright
  \item a collision, a blow, stumbles, injured, faints
  \item consciousness, lifts up, a litter, an emergency, leans on
  \item a forest fire, blazes up, a spark, dense, spreads
  \item fuel, a water jar, a stream of water, puts out a fire, dies down
  \item damage, destroys, broken, a citizen, what is left
  \item announces, danger, careful, keeps away, a misfortune
  \item an earthquake, shakes, slightly, cracks open, a pit
  \item a cloud, a whirlwind, whirls, flashes, uproots
\end{itemize}

\begin{tcolorbox}[breakable,colback=teal!4,colframe=teal!35!black,title={Before you move on}]
A collision? (\textbf{\sk{संघट्टः}}, \emph{saṅghaṭṭaḥ}.) A blow? (\textbf{\sk{घातः}}, \emph{ghātaḥ}.) Stumbles? (\textbf{\sk{स्खलति}}, \emph{skhalati}.) Injured? (\textbf{\sk{आहतः}}, \emph{āhataḥ}.) Faints? (\textbf{\sk{मूर्छति}}, \emph{mūrchati}.) Consciousness? (\textbf{\sk{चेतना}}, \emph{cetanā}.) Lifts up? (\textbf{\sk{उत्थापयति}}, \emph{utthāpayati}.) A litter? (\textbf{\sk{शिबिका}}, \emph{śibikā}.) An emergency? (\textbf{\sk{आपत्कालः}}, \emph{āpatkālaḥ}.) Leans on? (\textbf{\sk{अवलम्बते}}, \emph{avalambate}.) A forest fire? (\textbf{\sk{दावाग्निः}}, \emph{dāvāgniḥ}.) Blazes up? (\textbf{\sk{प्रज्वलति}}, \emph{prajvalati}.) A spark? (\textbf{\sk{स्फुलिंगः}}, \emph{sphuliṅgaḥ}.) Dense? (\textbf{\sk{गाढः}}, \emph{gāḍhaḥ}.) Spreads? (\textbf{\sk{प्रसरति}}, \emph{prasarati}.) Fuel? (\textbf{\sk{इन्धनम्}}, \emph{indhanam}.) A water jar? (\textbf{\sk{घटः}}, \emph{ghaṭaḥ}.) A stream of water? (\textbf{\sk{धारा}}, \emph{dhārā}.) Puts out a fire? (\textbf{\sk{निर्वापयति}}, \emph{nirvāpayati}.) Dies down? (\textbf{\sk{शाम्यति}}, \emph{śāmyati}.) Damage? (\textbf{\sk{क्षतिः}}, \emph{kṣatiḥ}.) Destroys? (\textbf{\sk{नाशयति}}, \emph{nāśayati}.) Broken? (\textbf{\sk{भग्नः}}, \emph{bhagnaḥ}.) A citizen? (\textbf{\sk{पौरः}}, \emph{pauraḥ}.) What is left? (\textbf{\sk{शेषः}}, \emph{śeṣaḥ}.) Announces? (\textbf{\sk{घोषयति}}, \emph{ghoṣayati}.) Danger? (\textbf{\sk{संकटम्}}, \emph{saṅkaṭam}.) Careful? (\textbf{\sk{सावधानः}}, \emph{sāvadhānaḥ}.) Keeps away? (\textbf{\sk{निवारयति}}, \emph{nivārayati}.) A misfortune? (\textbf{\sk{आपत्तिः}}, \emph{āpattiḥ}.) An earthquake? (\textbf{\sk{भूकम्पः}}, \emph{bhūkampaḥ}.) Shakes? (\textbf{\sk{कम्पते}}, \emph{kampate}.) Slightly? (\textbf{\sk{ईषत्}}, \emph{īṣat}.) Cracks open? (\textbf{\sk{स्फुटति}}, \emph{sphuṭati}.) A pit? (\textbf{\sk{गर्तः}}, \emph{gartaḥ}.) A cloud? (\textbf{\sk{मेघः}}, \emph{meghaḥ}.) A whirlwind? (\textbf{\sk{चक्रवातः}}, \emph{cakravātaḥ}.) Whirls? (\textbf{\sk{घूर्णते}}, \emph{ghūrṇate}.) Flashes? (\textbf{\sk{स्फुरति}}, \emph{sphurati}.) Uproots? (\textbf{\sk{उन्मूलयति}}, \emph{unmūlayati}.)
\end{tcolorbox}
